% =======================================================
% 2. QUALITA DI PROCESSO
% =======================================================
\section{Qualità di Processo}
Il gruppo adotta lo standard \textbf{ISO/IEC 12207} per la gestione dei processi. La qualità è perseguita attraverso il monitoraggio costante di indicatori quantitativi calcolati al termine di ogni Sprint\textsubscript{G}.

\subsection{Processi primari}
\subsubsection{Fornitura}
Questo processo comprende le attività volte a determinare le procedure per la pianificazione e il controllo del progetto. Il Responsabile\textsubscript{G} monitora l'avanzamento dei lavori e il consumo del budget\textsubscript{G} per garantire trasparenza alla proponente Zucchetti.

\subsubsection{Sviluppo}
Include le attività di analisi, progettazione e codifica. In questa fase di RTB, il focus è sulla stabilità dei requisiti estratti dal capitolato C6, al fine di minimizzare revisioni costose nelle fasi successive.

\subsection{Processi di supporto}
\subsubsection{Documentazione}
Assicura che ogni documento prodotto sia leggibile e conforme agli standard fissati nelle Norme di Progetto. La qualità è misurata tramite l'indice di leggibilità e l'assenza di errori formali.

\subsubsection{Gestione configurazione}
Garantisce l'integrità dei prodotti di lavoro tramite il versionamento su repository\textsubscript{G} Git. Ogni modifica deve essere tracciabile e associata a una specifica issue\textsubscript{G}.

\subsubsection{Verifica}
Accerta che ogni attività non contenga errori tramite un processo di controllo incrociato tra i membri del gruppo. Per i documenti, la verifica manuale si focalizza su sintassi, semantica e aderenza alle norme.

\subsubsection{Validazione}
Fornisce la conferma che il prodotto finale sia conforme alle attese dell'utente e della proponente. In questa fase si concentra sulla verifica che i requisiti individuati rispondano agli obiettivi del "Second Brain".

\subsection{Processi organizzativi}
\subsubsection{Gestione processo}
Si occupa del coordinamento del team durante i cicli di Sprint. Mira a bilanciare il carico di lavoro tra i componenti del gruppo Synergo Unit per evitare colli di bottiglia\textsubscript{G}.

\subsubsection{Miglioramento continuo}
Attraverso le retrospettive di fine Sprint, il gruppo analizza le inefficienze e definisce azioni correttive immediate per ottimizzare le metodologie di lavoro.

\subsection{Metriche per processi}
Le metriche di processo (MPC) vengono calcolate a ogni fine Sprint e riportate nel Cruscotto delle Metriche.

\begin{table}[H]
    \centering
    \renewcommand{\arraystretch}{1.3}
    \begin{tabularx}{\textwidth}{|l|X|c|c|}
        \hline
        \rowcolor{primarycolor!10}
        \textbf{Codice} & \textbf{Nome Metrica} & \textbf{Val. Accettabile} & \textbf{Val. Ottimale} \\
        \hline
        MPC01 & \textbf{Schedule Performance Index\textsubscript{G} (SPI)}: efficienza nel rispettare i tempi pianificati. & $\ge 0.9$ & $1.0$ \\
        \hline
        MPC02 & \textbf{Cost Performance Index\textsubscript{G} (CPI)}: efficienza nell'uso del budget allocato. & $\ge 0.9$ & $1.0$ \\
        \hline
        MPC03 & \textbf{Requirements Stability Index\textsubscript{G} (RSI)}: variazione dei requisiti nell'Analisi. & $\ge 80\%$ & $100\%$ \\
        \hline
        MPC04 & \textbf{Time Efficiency}: percentuale di ore produttive sul totale ore spese. & $\ge 75\%$ & $\ge 90\%$ \\
        \hline
        MPC05 & \textbf{Task\textsubscript{G} Completion Rate}: percentuale di issue chiuse rispetto a quelle pianificate. & $\ge 80\%$ & $100\%$ \\
        \hline
        MPC06 & \textbf{PR Resolution Time}: tempo medio per revisione\textsubscript{G} e merge\textsubscript{G} di una Pull Request\textsubscript{G}. & $\le 3 \text{ giorni}$ & $\le 1 \text{ giorno}$ \\
        \hline
        MPC07 & \textbf{Review Coverage}: percentuale di documenti/file verificati da un membro terzo. & $100\%$ & $100\%$ \\
        \hline
        MPC08 & \textbf{Quality Metrics Compliance Rate}: percentuale di metriche i cui valori sono almeno accetabili. & $\ge 80\%$ & $\ge 95\%$ \\
        \hline
    \end{tabularx}
    \caption{Metriche di Qualità di Processo}
\end{table}